% ================================================================
%  第 2 章 巡游一部作品——或孩子与母亲 (原书 p.31-53, 边码 P1-P37)
%  源文: _work/src1/task-ch02a.md
%  说明: 章内 Contents 目录不排印
% ================================================================
\chapter{巡游一部作品——或孩子与母亲}

\noindent 1986年1月\par

\section{(1) 事物的魔法}

% ---------------- p.31 ----------------
\origpage{P1} 小时候,我很喜欢上学。读和写、算术、唱歌(他拉一把小提琴给我们伴奏),还有史前人类和火的发现,都是同一位老师教我们。我不记得那时候我们在学校里曾经觉得无聊。他掌握着数字的魔法,也掌握着词语、符号和声音的魔法。还有韵脚的魔法,在歌谣里,在小小的诗里。韵脚里似乎藏着某种超出词语之外的奥秘。这情形一直持续到有一天
% ---------------- p.32 ----------------
有人向我解释,那不过是个再简单不过的“窍门”:所谓押韵,无非是让两个前后相接的说话片段以同一个音节收尾,于是它们就像中了魔法一样,变成了诗句。这可是一场启示!在家里,我身边有人接应,一连几个星期、几个月,我都玩着作诗。有一阵子,我说话只押韵了。幸好这毛病过去了。可即便今天,偶尔我还会写写诗——只是不大再去找韵脚了,除非它自己跑来。

另有一回,一个比我大的伙伴,已经在读中学了,教了我负数。那是另一种很好玩的游戏,但很快玩腻了。还有填字游戏——我花了好几天、好几个星期来编,越编越交错复杂。这游戏里结合着形式的魔法,以及符号和词语的魔法。但这份热情也离开了我,似乎没留下什么痕迹。

在中学里,先是在德国读第一年,后来在法国,我是个不错的学生,但算不上“出色的学生”。凡是我最感兴趣的,我不计工本地投入;不太感兴趣的,就往往被撂在一边,也不太在乎那位“老师”怎么评价。在法国上中学的第一年,1940年,我和母亲一起被关在芒德附近的里厄克罗集中营。那是战争年代,我们是外国人——用当时的话说,是“不受欢迎的人”。但集中营的管理当局对营里的孩子睁一只眼闭一只眼,尽管他们同样是“不受欢迎的人”。我们进出多少随自己的便。我年纪最大,也是唯一去上中学的,学校在四五公里外,不管下雪还是刮风,脚上穿着总是进水的破鞋。\origpage{P2} 我还记得第一次“数学考试”,老师给了我一个坏分数,为的是证明“三角形全等的三种情形”之一。我的证明不是书上那个,而他是虔诚地照书办的。可我心里很清楚,我的证明并不比书上那个更有说服力,也不更没说服力——我遵循的正是书里的精神,用的还是那套老掉牙的“把某个图形这样挪到另一个图形上去”。显然,教我的这个人觉得自己没有能力凭自己的眼光来判断(在这里,就是判断一个推理是否成立)。他必须求助于一种权威,此处即一本书的权威。这些做派一定深深触动了我,我才会记得这桩小事。此后直到今天,我有很多机会看到:这样的做派绝不是例外,而是几乎普遍的规则。这个话题有很多可说的——在《收获与播种》里,我不止一次以这种或那种形式触及它。可直到今天,不管我愿不愿意,每次再撞上它,我都感到手足无措……

战争最后几年,母亲仍被关在集中营里,我则待在“瑞士救援会”办的一所儿童之家,收容难民儿童,在利尼翁河畔勒尚邦。我们大多是犹太人;每当接到通知(由当地警察来通知)说盖世太保要来搜捕,我们就到树林里躲上一两夜,三三两两一小群,并不太意识到这真的关系到我们的性命。那一带塞文山区藏满了犹太人,许多人靠当地居民的互助活了下来。

在“塞文中学”(我在那里受教育)最让我吃惊的,是我的同学们对所学的东西兴趣之少。至于我,学年一开始就把课本狼吞虎咽地读完,心想这一回总算要学到真正有趣的东西了;而这一年余下的时间,我就尽量自己打发,与此同时,预定的教学大纲一个学期接一个学期,无情地照本宣科。不过我们的老师都很好。博物学老师弗里德尔先生,人品与学识都极出色。可他不会“镇压”,被人起哄到死,以致到了学年末,课已经没法再上下去了,他那无力的声音
% ---------------- p.33 ----------------
被一片喧闹淹没。说不定正因为这个,我才没成为生物学家!

\origpage{P3} 我相当多的时间——甚至上课的时候(嘘……)——都在做数学题。很快,书上的题就不够我做了。也许是因为它们久而久之彼此太相像了;但我相信,尤其是因为它们来得太像从天而降,就这么一个接一个排着队,不说自己从哪儿来、往哪儿去。那是书上的问题,不是我的问题。其实真正自然的问题并不缺。比如,已知三角形三边之长 $a$、$b$、$c$,这个三角形就确定了(不计其位置),那么就该有一个显式的“公式”,例如把三角形的面积表示成 $a$、$b$、$c$ 的函数。四面体也一样:已知六条棱的长,体积是多少?这一回我想我大概费了不少劲,但硬啃下去,终究还是做出来了。反正,一件事“抓住”了我,我就不数在上面花了多少小时、多少天,哪怕把别的全都忘了!(如今也还是这样……)

我们的数学书里最让我不满意的,是对长度(一条曲线的)、面积(一个曲面的)、体积(一个立体的)这些概念,完全没有认真的定义。我向自己保证,一旦有闲,就要补上这个空缺。1945到1948年间,我在蒙彼利埃大学念书,把最精壮的精力都花在了这上面。学院的课不合我的意。虽然从未对自己明说,我大概觉得老师们不过是在重复他们的书,就像芒德中学那位第一位数学老师一样。所以我也只是隔很久才去一趟学院,为的是了解那个老掉牙的“教学大纲”讲到了哪里。就那大纲而言,书是够用的;但同样清楚的是,它们根本不回答我给自己提的问题。说实话,它们连看都没看见那些问题,就像我的中学课本也没看见一样。只要它们随手给出计算长度、面积、体积的配方,靠单重、双重、三重积分(超过三维的情形则谨慎地绕开……),那么为它们给出一个内蕴的定义这回事,似乎就根本不成为问题——对我的老师们如此,对教科书的作者们也如此。

\origpage{P4} 以我当时有限的经验看,完全可能显得我是世界上唯一对数学问题怀有好奇心的人。无论如何,在那些年里,这就是我未曾说出口的信念;那些年我处于完全的智识孤独之中,而这份孤独并不使我难受\fnote{1}{}{1945到1948年间,我和母亲住在蒙彼利埃十来公里外的一个小村子里,迈拉尔格(属旺达尔格乡),四野是葡萄园中的一处僻静地。(父亲1942年在奥斯维辛失踪。)我们靠我那份微薄的学生助学金过活,过得很拮据。为了收支相抵,我每年去摘葡萄,摘完葡萄再做些“拾穗酒”,勉强脱手(据说这违反了当时的法律……)。此外还有一处园子,从不用侍弄,却大量供给我们无花果、菠菜,甚至(到最后)还有番茄——那是一位好心的邻居种下的,种在一片绚丽的罂粟正中间。日子过得美——只是有时紧得刚好差那么一点,比如要换一副眼镜架,或者一双穿到露出麻线的鞋。幸而母亲在集中营里待了那么久,身体虚弱多病,可以享受免费医疗。否则我们永远付不起看医生的钱……}。说实话,我想那时候我从未深究过:自己是否真是世上唯一可能对我所做的事感兴趣的人。我的精力已充分投入那场自己给自己出的挑战:发展一套能让我完全满意的理论。

我丝毫没有怀疑自己终能做到、终能找到事物的最终答案——只要我肯花力气去细究它们,把它们对我说的随说随记到纸上。比如说,体积的直觉是无可置疑的。它只能是一个现实的映照,这现实眼下还难以捉摸,但完全可靠。要把握的就是这个现实,如此而已——也许有点像有一天“韵脚”那个神奇的现实被把握住、“被理解”那样。

% ---------------- p.34 ----------------
动手的时候,我十七岁,刚出中学,我以为这不过是几个星期的事。结果我在上面待了三年。我还想方设法,硬是在大学二年级期末考砸了一门——球面三角(在“深入天文学”选修方向里,原文如此),因为一个愚蠢的数值计算错误。(得承认,一离开中学,我的计算就一直不怎么行……)正因如此,我才不得不在蒙彼利埃再待第三年,在那里读完学士学位,而不能马上去巴黎——据人向我保证,那是唯一有机会遇到懂行之人的地方,懂数学里什么才算重要。给我通风报信的苏拉先生还向我保证:数学里最后剩下的一些问题,二三十年前已由一位名叫勒贝格的人解决了。他恰好发展了一套测度与积分理论,(说真的,巧得可笑!)给数学画上了句号。

\origpage{P5} 苏拉先生教我们“微积分”,是个和善的人,对我一片好意。我并不认为他因此说服了我。我心里大概早已预感到:数学在广度与深度上都是无边无际的。大海有“句号”吗?总之,我从未动过念头去把苏拉先生提到的勒贝格那本书找来,而他自己大概也从未把这本书拿在手里过。在我心里,一本书里可能装的东西,和我为了满足自己对那些令我不解之事的求知欲、照自己的方式所做的工作,毫无共同之处。

\section{(2) 独处的重要}

一两年后,我终于在巴黎与数学界有了接触,最终在那里得知(在许许多多别的之外):我在自己角落里用凑合的手段做出来的工作,(几乎)就是“人人”都熟知的东西,叫做勒贝格测度与积分理论。在我向其谈起这项工作(甚至给他们看一份手稿)的两三位年长者眼里,这有点像我白白浪费了时间,重做了一遍“已经知道的东西”。我也不记得自己有过失望。那时候,为自己所做的工作去获取什么“信誉”、哪怕只是一句赞许或别人一点兴趣,这个念头大概还离我的头脑很远。何况我的精力已经够忙的了:要熟悉一个完全不同的圈子,尤其是要学会巴黎人眼中数学家的 ABC\fnote{2}{}{我在《收获与播种》第一部分(ReS I)“受欢迎的异乡人”一节(第9号)里,对这段有点粗粝的过渡时期作了简短的叙述。}。

然而,现在回想这三年,我意识到它们绝非白费。在不知不觉中,我在孤独里学到了数学家这一行当最本质的东西——那是任何老师都无法真正教给人的。我从未需要对自己说明,也从未遇到一个可以分享我求知之渴的人,可我却知道,用俗话说,“打骨子里”知道:我是一个数学家——一个在完全意义上“做”数学的人,就像人们“做”爱那样。数学对我来说,成了一位永远迎候我欲望的情人。这些孤独的岁月奠定了一份从未动摇过的信赖的基础——无论是二十岁到巴黎时发现自己的无知有多么深、要学的东西有多么浩瀚;还是二十多年后,我永远离开数学界时那些动荡的插曲;还是最近这些年,由我昔日最亲近的伙伴们操办的那场针对我本人和我的作品的、相当疯狂的“葬礼”
% ---------------- p.35 ----------------
(提前举行,而且办得干净利落)……\origpage{P6}

换一种说法:在这些关键的年月里,我学会了独处\fnote{3}{}{这个说法有些不确切。我从来不需要“学会独处”,原因很简单:在我的童年里,我从未丢掉出生时就在我身上、也就在每个人身上的那种天生能力。但这三年孤独的工作——在其中我得以按自己那套自发的严格要求的标准,给自己量了尺寸——却在我身上确认并重新奠下了一份信赖与安然的确据,这一次是在我与数学工作的关系上;这份确据与那些立法的共识和风尚毫无关系。我在《根与孤独》一注(ReS IV,第1713号,尤其第1080页)里有机会再次提到它。}。我的意思是:凭自己的眼光去接近我想认识的事物,而不是信赖那些明说的或默认的观念与共识——它们来自我自认为属于的某个或大或小的群体,或由于任何别的理由而被赋予对我的权威。中学也好,大学也好,那些无声的共识曾告诉我:对“体积”这个概念本身,没有必要提出问题,它被说成“人所共知”“显而易见”“没有问题”。我越过了它,当作理所当然——正如几十年前勒贝格想必也曾越过它。正是在这种“越过”的行动里,在总之做自己、而不只是做那些立法共识的表达式里,在不被它们为我们划定的那个强制性圈子关在里面——首先是在这种孤独的行动里,才有“创造”。其余的一切都是额外来的。

后来,在接纳我的这个数学界里,我有机会遇到许多人,无论是年长者,还是与我年纪相仿的年轻人,他们显然比我聪明得多、也“有天赋”得多。我欣赏他们学新概念时那种轻松,像玩一样,拿来耍弄,仿佛从摇篮里就认识它们——而我却觉得自己笨重迟钝,像一只鼹鼠,在由那些(人们向我保证)我理当学会的东西堆成的一座不成形的山里面艰难地开路,而那些东西的来龙去脉我根本抓不住。事实上,我一点也没有那种出色学生的样子:轻松通过名校会考,三下两下就消化掉那些吓人的教学大纲。

\origpage{P7} 何况,我那些更出色的同学大多成了有本事、有名望的数学家。然而,隔了三十或三十五年回头看,我发现他们并没有在我们这个时代的数学上留下真正深刻的印记。他们在已经现成的环境里做了些事情,有时是漂亮的事情,而那环境他们从没想过要去触动。他们不知不觉地成了那些看不见而专横的圈子的囚徒——那些圈子在某个圈子、某个时代里界定出一个“宇宙”。要跨越它们,他们本应在自己身上重新找到出生时即有的那种能力,那种我也有的能力:独处的能力。

小孩子呢,独处毫无困难。他天生就是孤独的,虽然偶尔有人做伴他也不讨厌,而且知道到了吃奶的钟点就吵着要妈妈的“咪咪”。他也清楚地知道,不用谁来告诉他,那“咪咪”是给他的,而他会吃。可我们常常已经和心里的那个孩子失去了联系。我们不断地从最好的东西旁边走过,不屑于看它一眼……

如果说在《收获与播种》里我除了对自己说话,还对别的什么人说,那绝不是对一个“公众”说。我在那里是对你、对读我的你说话,把你当作一个人,一个孤独的人。我想说话的,是你身上那个懂得独处的人,那个孩子,而不是别人。那孩子常常很远,我很清楚。他什么没见过,而且早就见惯了。他不知躲到哪里去了,要走到他跟前,常常不容易。人家会赌咒说,他早就死了,或者干脆从未存在过——可我确信他就在某处,活得好好的。

我也知道我被听懂的标志是什么。那就是:超越一切文化与命运的差异
% ---------------- p.36 ----------------
,我所说的关于我本人、我的生活的话,在你那里得到回声与共鸣;当你在其中也重新找到你自己的生活、你自己对自己的经验,也许还是你此前未曾留意的一面。这不是什么“认同”,不是去认同与你相距遥远的什么东西或什么人。而也许是,你稍微地,通过我对自身生活的重新发现,也重新发现了你自己的生活,那离你最近的东西——顺着《收获与播种》一页一页,直到我今天此刻正在写下的这些页面。\origpage{P8}

\section{(3) 内在的冒险——或神话与见证}

首先,《收获与播种》是一场对我自己、对我的生活的反思。因此它也是一份见证,而且是从两方面来说的见证。它是对我过去的见证,反思的主要重量正落在这过去上。但同时它也是对最切近的当下的见证——对我此刻正在写作的见证,对《收获与播种》这些页面在一个个钟点、一个个夜晚和白天里诞生的此刻的见证。这些页面是一场关于我的生活的长久沉思的忠实见证,那沉思如它实际所经历的那样延续着(并且此刻仍在延续……)。

这些页面没有文学上的野心。它们是一份关于我自己的文献。我只在极狭窄的限度内才允许自己动它(比如偶尔作些文体上的修饰)\fnote{4}{}{因此,对错误(事实上的,或视角上的,等等)的可能的更正,并不是对初稿加以修补的机会,而是放在页底注里,或在日后对所述情境的“回访”中进行。}。若说有什么野心,那也只是:是真的。而这已经很多了。

此外,这份文献一点也不是“自传”。你在里面不会知道我的出生日期(它除了用来画一张星象图,恐怕没什么意思),也不会知道我母亲和父亲的名字、他们以什么为生,不会知道那位曾是我的妻子的人和其他在我生命中有过重要位置的女人的名字,也不会知道这些爱情中生下的孩子们的名字,以及他们各自把自己的生活过成了什么样。这并不是因为这些事情在我生命里不重要,也不是因为它们今天已不再重要。而是因为,这场对我的反思既已如此开始并延续下去,我从未在某一刻感到被促使去稍稍着手描述这些我在这里那里掠过的事情,更不会去认真地罗列名字和数字。从未有一刻,我会觉得这能给当时我所追索的“话题”增添什么。(而在前面那几页里,我却像不由自主似的,写进了也许比接下来一千页里还要多的、关于我生活的物质细节……)

如果你问我,这千页一路追索的“话题”究竟是什么,我会回答:是讲述、并由此发现我一生所是、且现在仍是的那场内在冒险。这份关于一场冒险的叙述—见证,同时在我刚才说的两个层面上展开。一个是对过去那场冒险的勘探,从它的根与起源一直追溯到我的童年。另一个是这“同一场”冒险的延续与更新,在我写《收获与播种》的一个个瞬间和一天天里,作为对来自外部世界的一声猛烈呼唤\fnote{5}{}{关于这声“猛烈呼唤”的细节,见《信》,特别是第3至8节。}的自发回应。\origpage{P9}

外部的事实为反思提供养料,但只在这样的限度内:它们激起并引发内在冒险的又一次跃起,或有助于照亮它。我的数学作品被埋葬、被洗劫这件事(后面会详谈),就是这样一次挑衅。它在我心里激起了强大的自我反应的全员出动,同时又向我揭示了那些深的、此前未被察觉的纽带,它们仍把我与从我这里生出的作品连在一起。

% ---------------- p.37 ----------------
\origpage{P10} 确实,我属于“数学尖子”这件事,并不必然成为你对我这段特殊冒险感兴趣的理由(更不是好理由)——我改了圈子、改了生活方式之后跟同行们闹出的麻烦,也不是理由。而且并不缺少同行,甚至朋友,他们觉得当众摊开自己的“心境”(用他们的话说)是最可笑不过的。要紧的是“成果”。“灵魂”呢,也就是说我们身上那个经历着这些“成果”的“生产”、以及它的各种余波(既在“生产者”的生活里,也在他的同类的生活里)的东西,却遭人轻视,甚至被公然嘲弄。这种态度自命为“谦虚”的表现;我在其中看到的却是一种逃避的征兆,一种古怪的失序,而且是我们呼吸的空气本身所助长的。我当然不是为那种被对自己的暗暗蔑视所击中的人写作的——那蔑视使他看不上我能给他的最好的东西。一种对真正构成他自己生活之物的蔑视,对构成我生活之物的蔑视:那些驱动着心灵的浅层与深层运动,粗糙的或细微的,那正是经历着经验、并对之作出反应的“灵魂”,它僵住或舒展,它退缩或在学习……

一场内在冒险的叙述,只能由经历它的人来做,别人不行。但即便这叙述只给自己看,它也很少不滑进建造神话的车辙里,而那神话的主角就是叙述者本人。这样的神话,并非诞生于一个民族、一种文化的创造想象,而是诞生于那个不敢承担一份谦卑现实的人的虚荣,他乐于用一件构造物、一件他心智的作品去替换它。但一份真实的叙述(如果有的话),关于一场真正被那样经历过的冒险,是珍贵之物。而这并不因为围绕叙述者的某种威望(不论有理没理),而只因为它存在,带着它的真之品质。这样的见证是珍贵的,无论它来自一位有名望甚至显赫的人物,还是来自一个没有前途、拖家带口的小职员,还是来自一个普通刑事犯。

这样一份叙述若对别人有什么功效,那首先是通过对另一个人经验的不加粉饰的见证,让他重新面对自己。或者(换个说法)也许是抹去他心中(哪怕只在读一次书那么长的时间里)那份蔑视——他借它来看待自己的冒险,以及那个既是乘客又是船长的“灵魂”……

\section{(4) 风俗画}

说到我作为数学家的过去,后来又(像违心地)发现我那场巨大“葬礼”的种种波折与内幕,我不由自主地——并非刻意——画出某个圈子、某个时代的画卷:一个某些价值正在瓦解的时代,那些价值曾赋予人的劳动以意义。这就是“风俗画”的一面,围绕一桩在“科学”编年史上大概是独一无二的“社会新闻”画成。我前面说的话,我想已经足够清楚:你在《收获与播种》里找不到一份关于某桩不寻常“事件”的“卷宗”,好让你很快了解情况。可是我有一位朋友,一直在找那份卷宗,他闭着眼睛走过,什么也没看见,几乎错过了构成《收获与播种》的实质与血肉的一切。

正如我在《信》里详细得多地解释的,“调查”(或“风俗画”)主要在第二部和第四部展开:《葬礼(一)——或中国皇帝的新衣》和《葬礼(三)——或四种运算》。一页一页地,我在那里执拗地把一大批(至少可以说)有汁水的事实一件一件挖出来,又尽量随挖随“安顿”。渐渐地,这些事实聚成一幅总体画面,它逐步从雾中显出
% ---------------- p.38 ----------------
,颜色越来越鲜,轮廓越来越清晰。在这些逐日写下的注里,刚出现的“原始事实”与个人的回忆、与心理的、哲学的、乃至(偶尔)数学的评论和思考,不可分割地混在一起。就是这样,我也无能为力!

以我做过的那项工作(它让我绷紧了一年多)为基础,用“调查结论”的体例编一份卷宗,应该是几个钟头或几天的额外工作,视有兴趣的读者的好奇与要求而定。我确实试过一回,要编那份著名的卷宗。那是我开始写一个本该叫《四种运算》的注的时候\fnote{6}{}{那条原计划的注最终炸开成了《收获与播种》第四部(同名《四种运算》)里的大约70条注,足足铺开四百来页。}。可后来,不行,一点办法也没有。我做不到!这确实不是我的表达方式,到了晚年更是如此。而且我现在认为,有《收获与播种》,我为“数学共同体”做的已经够多了,可以毫无悔意地把编那份该编的“卷宗”的事,留给我以外的人(如果我的同行中有谁觉得自己与此有关)。\origpage{P11}

\section{(5) 继承人与建造者}

现在该在这里说几句我的数学作品了——它在我生命里占据了一个重要的位置,而且(令我自己惊讶)至今仍占着。《收获与播种》里我不止一次回到这部作品——有时用每个人都能清楚明白的方式,有时则用多少带点技术性的词语\fnote{7}{}{除了关于我过去作品的数学概览,这里那里还有一些含新数学发展的段落。最长的一篇是 ReS IV 第171(ix)号注里的《五张照片(晶体与 D-模)》。}。后一类段落,在很大程度上会“越过”不止“门外汉”的头顶,连对其中所谈的数学已经不怎么“在行”的同行数学家也一样。你当然可以干脆跳过那些你觉得太“专业”的段落。你也可以把它们浏览一遍,也许顺带捕捉到数学事物世界那“神秘的美”(这是一位非数学家朋友来信中的说法)的一线反光——它们像一座座“不可企及的奇异小岛”,在反思的浩瀚流动的水面上浮现……

大多数数学家,我刚才说过,倾向于把自己圈在一个概念框架里、一个一劳永逸地定下来的“宇宙”里——基本上就是他们求学时发现已经“现成”的那个。他们像一座宽敞漂亮的、一切都安置好了的房子的继承人:有客厅、厨房、作坊,有成套的锅碗瓢盆和随手可用的工具,凭这些,天哪,足够做饭和做点手工了。这房子在几代人手里是怎样一点点盖起来的,某些工具(而不是别的……)是怎样、为什么被设计和打造出来的,为什么这间屋子这样布置、那间那样安排——这些问题是这些继承人从不会想到去问的。这就是“宇宙”,就是必须生活在其中的“给定之物”,如此而已!它显得很大(而且多数时候,离把它的所有房间都走遍还远着呢),同时又很熟悉,尤其是:不可变易。他们忙活,是为了维护和装点一份家产:修一张摇晃的柜子,给外墙抹灰,磨快一件工具,甚至有时候,对比较有事业心的来说,在作坊里从头做一件新家具。而他们全心投入时,这件家具也可能美极了,整座房子也显得更美了。\origpage{P12}

更罕见的是,他们中有人会想到去改动一件备用的工具,或者甚至在需求的反复催促下,去设想并造出一件新的。这样做的时候
% ---------------- p.39 ----------------
,他几乎要连声道歉,因为他觉得自己这样有点像冒犯了家族传统应有的虔敬,觉得自己用一桩古怪的创新把它冲撞了。

房子里大多数房间的窗和百叶窗都仔细关着——大概是怕有一股别处来的风灌进来。而当漂亮的新家具这里一件、那里一件,还不算子嗣,开始把变得狭窄的房间塞满,一直侵占到走廊时,这些继承人里没有一个愿意意识到:他那熟悉而舒适的宇宙开始有点紧得转不开身了。与其承认这一点,他们各人宁愿勉勉强强地侧身挤过去、卡在那里:有人挤在路易十五式碗柜和藤编摇椅之间,有人挤在流鼻涕的娃娃和埃及石棺之间,最后还有一位,在无计可施之下,尽其所能地爬上一堆杂七杂八、摇摇欲坠的椅子和长凳……

我刚画的这幅小画,并非数学界所特有。它说明的是一些根深蒂固、由来已久的制约,在所有圈子和人类活动的一切领域里都能遇到,而且(据我所知)在所有社会、所有时代都是如此。我此前已有机会提及,我也绝不声称自己可以免于它。我的见证将表明,实情恰恰相反。只不过,在智力创造活动这个相对有限的层面上,我受这种制约的触动相当少\fnote{8}{}{我认为主要原因在于某种有利的氛围,它环绕着我的童年,直到我五岁。参见《无辜》一注(ReS III,第107号)。}——它可以叫做“文化盲”:看不见(也无法活动于)周围文化所固定的那个“宇宙”之外。

至于我,我觉得自己属于那样一脉数学家:他们自发的天职与快乐,是不断地建造新房子\fnote{9}{}{这个“待建之屋”的原型意象,第一次浮现并被表述出来,是在《阴,仆役,与新主子们》一注中(ReS III,第135号)。}。一路之上\origpage{P13},他们也忍不住随时发明并打造一切必需的工具、器具、家具和仪器,既为了从地基到屋脊把房子建起来,也为了给未来的厨房和作坊备足东西,把房子布置得能住、能舒服。然而,一旦连最后一条檐沟、最后一只凳子都安放好了,这位工匠却很少在这些地方久留——这里的每一块石头、每一根椽子,都带着加工并安放它的那只手的痕迹。他的位置不在那些现成宇宙的安宁里,无论它们多么宜人、多么和谐——无论是他亲手布置的,还是他的前人们布置的。别的任务已经在召唤他去新的工地:在某种需要的专横推动下,那需要也许只有他一个人清楚地感到,或者(更常见的是)他预感到某种只有他一个人预感到的需要。他的位置在露天里。他是风的朋友,不怕独自干活,一干几个月、几年,必要时干一辈子,除非有可喜的接替者来搭把手。他当然和所有人一样只有两只手——但那是两只随时能猜到自己该做什么的手,既不嫌最粗笨的活,也不嫌最精细的活,而且从不厌倦去认识、反复认识那些无数次召唤它们去认识的事物。两只手也许太少,因为世界是无穷的。它们永远也穷尽不了它!然而,两只手又很多……

我这个人历史不行,但若要我举出这一脉数学家的名字,我自然而然想到伽罗瓦和黎曼(上世纪),以及希尔伯特(本世纪初)。若要在当初接纳我进入数学界的年长者中找一位代表\fnote{10}{}{我在《受欢迎的异乡人》一节(ReS I,第9号)里谈到这些开端。},我最先想到的是让·勒雷(Jean Leray)的名字,尽管我和他的接触
% ---------------- p.40 ----------------
始终极为零星\fnote{11}{}{不过这并不妨碍我(继 H. 嘉当与 J.-P. 塞尔之后)成为勒雷引入的某个重大革新概念——层——的主要使用者和推广者之一,而层贯穿我整个几何学家的工作,是其中一件本质的工具。也正是它给了我钥匙,把(拓扑)空间概念扩展为 topos 概念,后面会谈到。\par 此外,在我看来,勒雷与我勾画的“建造者”肖像有所不同:他似乎并不倾向于“从地基到屋脊建造房子”。不如说,他忍不住在无人想得到的地方开出了广阔的地基,却把完成它们、在上面建造、以及房子建成后住进去(哪怕只住一阵)的事留给别人……}。

我刚才大致勾画了两幅肖像:一位是“恋家”的数学家,满足于维护和装点一份遗产;一位是建造者—拓荒者\fnote{12}{}{我刚刚偷偷地、“从侧面”把两个带有雄性回响的修饰语(即“建造者”和“拓荒者”)并置在一起,而它们表达的却是发现冲动相当不同的方面,其本性比这些名字所能唤起的更精细。这将在这次巡游—反思的后文、“发现母亲——或两个山坡”(第17号)那一步里显现。},他忍不住不断越过那些界定一个宇宙的“看不见而专横的圈子”\fnote{13}{}{而且,他也就这样在无意中给这个旧宇宙(不是为他自己,至少是为那些不如他好动的同类)指定了新的界限,划在新的圈子里——这些圈子当然更广,但和它们所取代的那些一样看不见、一样专横。}。也可以用两个有点一刀切、但颇有意味的名字称他们:“守成者”与“革新者”。两者在同一个延续于世代、世纪、千年的集体冒险中,都有存在的理由和要扮演的角色。在一门科学或一门艺术的繁盛期,这两种气质之间没有对立,也没有敌意\fnote{14}{}{尤其在数学界,在我亲身见证的时期(1948–1969),情况就是如此,当时我自己也是这个世界的一部分。1970年我离开之后,似乎出现了一种大范围的反应,一种对“想法”总体的“蔑视的共识”,尤其针对我引入的那些重大革新想法。}。它们不同,而且互相补充,就像面团与酵母互相补充。\origpage{P14}

在这两个极端类型之间(它们在本性上丝毫不相反对),当然有一整套中间气质。某位“恋家者”从不想离开一处熟悉的住所,更不会去揽下到天知道什么地方另建一座的活计;然而当住处真的开始变窄时,他会毫不犹豫地拿起抹刀,收拾一间地窖或阁楼,加高一层,必要时还在墙边接出一间规模不大的新附属房\fnote{15}{}{我的大多数“前辈”(例如《一笔情愿的债》导论第10节里谈到的那些)都属于这种中间气质。我特别想到亨利·嘉当、克劳德·谢瓦莱、安德烈·韦伊、让–皮埃尔·塞尔、洛朗·施瓦茨。除也许韦伊之外,他们对我眼看要登上的那些孤独冒险,都投以“善意的目光”,没有“暗暗的担忧或责备”。}。他骨子里不是建造者,但常常会带着善意的眼光——至少是不带暗暗的担忧与责备——看着另一位曾与他同住一座房子的人,如今
% ---------------- p.41 ----------------
在某个穷乡僻壤辛苦地聚拢梁木和石头,一副仿佛已经看见宫殿的神气……\origpage{P15}

\section{(6) 视角与眼光}

但还是回到我自己和我的作品上来。

如果说我在数学家的技艺上有所出众,那与其说靠解题的灵巧与坚持(解前人传下的问题),不如说靠我身上这种天然的倾向:它驱使我看见一些显然关键、却无人看见的问题,或者提取出那些缺失的“好概念”(往往在新概念出现之前,谁也没意识到它们缺失),以及那些谁也没想到的“好陈述”。常常,概念与陈述配合得如此完美,使我心里不可能怀疑它们是对的(顶多需要些修补)——于是往往,当这不是要发表的“部件活”时,我就省得再往前走,省得花时间去把证明弄妥;而那证明,一旦陈述和它的语境看清了,多半只是个“手艺”问题,甚至可以说是例行公事。值得注意的事情数不清,不可能把每一个的召唤都跟到底!尽管如此,在我已写出并发表的作品里,正式证明过的命题和定理数以千计,我想我可以说,除了极少的例外,它们全都进入了共同的遗产,被普遍当作“已知”,在数学各处被日常使用。

但我的独特天赋把我带向的,比发现新问题、新概念、新陈述更进一步的,是发现富有成果的视角——它不断引导我引入并或多或少地发展一些全新的主题。我觉得,这就是我对我那个时代的数学所贡献的最本质的东西。说实话,我刚才说的那无数问题、概念、陈述,在我看来只有在这样一种“视角”的光照下才有意义——更确切地说,它们是从中自发地、带着显然性之力生出来的;就像黑夜里亮起的一道光(哪怕是散漫的光),仿佛从虚无中生出它突然向我们揭示的那些或模糊或清晰的轮廓。没有这道把它们束成一束的光,那十个、一百个、一千个问题、概念、陈述,就会显得像一堆杂七杂八、没有形状的“心智小玩意儿”,彼此孤立——而不是一个整体的各个部分;那整体也许仍不可见,仍躲在夜的褶皱里,却已被清楚地预感到。

富有成果的视角,就是那样一种视角:它把那些无人感到的灼热问题,揭示为同一个包容它们、赋予它们意义的整体的活的各部分;并且(也许是对这些问题的回答)把那些如此自然、却无人想到去提取的概念,以及那些仿佛顺理成章、却在激发它们的问题和使之得以表述的概念尚未出现之前谁也不敢提出的陈述,一并这样揭示出来。在我们这门技艺里\fnote{16}{}{毫无疑问,情况不只在我们这门“技艺”里如此,而且(在我看来)在一切发现的工作里都是如此,至少当这工作处于智识认识的层面时。},比所谓“关键定理”更有力的发现工具,是富有成果的视角——更确切地说,它们不是工具,而是研究者本人的眼睛,那双热切地想要认识数学事物之本性的眼睛。

这样,富有成果的视角不外是那只“眼睛”:它同时让我们发现,又让我们在被发现者的多样性中认出统一。而这统一,真正是把这些多样之物联系起来并赋予生命的生活本身与气息。

但正如其名所示,一个“视角”本身仍是局部的。它向我们揭示一处风景或全景的一个侧面,而此外还有许多同样有效、同样“真实”的侧面。只有当对同一现实的互补视角彼此结合起来,当我们的“眼睛”增多,目光才更深入地进入对事物的认识。我们想认识的现实越丰富、越复杂,拥有好几只“眼睛”\fnote{17}{}{每个视角都会带来发展出一套表达它、并专属于它的语言。拥有好几只“眼睛”或好几个“视角”来把握一个情境,也等于(至少在数学里)拥有好几种不同的语言来圈定它。}以便在它的全部广度与全部精微中把握它,也就越重要。

而有时候,一束汇聚在同一片广阔风景上的视角,凭着
% ---------------- p.42 ----------------
在我们身上那种能在多中把握一的能力,使一件新事物有了身体;这件新事物超出每一个局部透视,正如一个活物超出它的各个肢体和器官。这件新事物,可以叫做一种眼光。眼光把体现它的那些已知视角统一起来,并向我们揭示另一些此前被忽略的视角,正如富有成果的视角让人发现并把握众多新问题、新概念、新陈述,把它们当作同一整体的部分。\origpage{P16}

换个说法:眼光对于它仿佛从中生出、并由它统一起来的那些视角,就像白昼明亮温暖的光对于太阳光谱的不同成分。一种广阔而深沉的眼光,像一处取之不竭的源泉,注定要去启发和照亮的工作,不仅是在某一天生出它、并使自己成为它的仆人的那个人的工作,也是世世代代人的工作——他们也许(像他本人一样)被它让我们隐约看见的那些遥远边界所迷住……\origpage{P17}

\section{(7) “大想法”——或树木与森林}

我数学活动所谓“多产”的时期,也就是有正式发表物作证的时期,在1950到1969年之间,共二十年。而在二十五年里,从1945年(我十七岁)到1969年(我将满四十二岁),我几乎把全部精力投进了数学研究。这投入无疑过分。我为此付出的是长期的精神停滞、逐步的“变厚”,在《收获与播种》的篇幅里我不止一次有机会谈到。然而,在一项纯智力活动的有限领域内,凭一种局限于数学事物世界的眼光的绽放与成熟,那些年仍是创造力旺盛的年份。

在我生命的这段漫长时期,几乎全部时间和精力都用在所谓“部件活”上:那种细致的打造、装配和磨合工作,是从头建造一座座房子所必需的,而一个内心的声音(或一个魔鬼……)命令我建造它们,并按工作进展随时向我口授一位总施工的方案。我被“手艺”的活计占住:时而当石匠,时而当砌工、木匠,甚至当铅管工、细木工和家具匠——我很少偷闲把那份除我之外无人可见的总图(后来才显出它的样子……)用白纸黑字记下来,那总图在日、月、年之中以梦游者般的稳当引导着我的手\fnote{18}{}{“梦游者”这个意象,是柯斯特勒(Koestler)那本出色的书《梦游者》的标题给我的(Calmann-Lévy 版),那本书是一篇“关于宇宙观念史的研究”,从科学思想的开端直到牛顿。这段历史中令柯斯特勒深受触动、并由他特别指出的一点是:我们的世界认识常常从某一点走到另一点,那另一点(从逻辑上、事后看来)显得近在咫尺,而道路却要经过最离奇的迂回,仿佛在挑战健全的理性;然而,尽管有这千百道似乎注定要让人永远迷路的迂回,出发去寻找宇宙之“钥匙”的人们,却带着“梦游者般的稳当”,不由自主地、往往甚至连自己也没有察觉,碰上了另一些他们远未预料到的“钥匙”,而那些竟然正是“对的”。\par 据我在周围所能观察到的,在数学发现的层面上,这些发现之路上的惊人迂回,是某些大格局研究者的事,但绝不是所有研究者的事。这也许是因为,两三个世纪以来,自然科学的研究,尤其是数学研究,已从与某一文化、某一时代相联系的宗教或形而上学强制前提中解放出来——那些前提曾是对宇宙的“科学”理解之展开(无论好坏)特别强大的制动。然而确实,数学中某些最根本、最显然的观念与概念(如位移、群、零这个数、文字运算、空间中一点的坐标、集合的概念,或拓扑“形式”的概念,更不必说负数与复数),都经过几千年才出现。这些都是那种根深蒂固、深深植入心灵之中的“阻塞”的明白迹象——它抵制全新的观念,即便这些观念天真地简单,似乎以显然之力自行强加于人,而且几代、乃至几千年都是如此……\par{}回到我自己的工作,我觉得其中的“失手”(也许比多数同行更多)只限于细节,通常很快就由我自己发现。那不过是纯属“局部”的“途中意外”,并不严重影响关于所考察情境的基本直觉的可靠性。相反,在观念与重大指导性直觉的层面上,我的作品似乎没有任何“失误”,尽管这听起来难以置信。每一刻,即便不能把握一条进路的最终归宿(它多半仍藏在目光之外),也至少能把握那些最富成果、引我直达本质事物的方向——正是这种从未失灵的可靠性,使柯斯特勒的“梦游者”形象重新浮上我的心头。}。应该说,我乐于在“部件活”里倾注一份充满爱意的细致,这活并不令我讨厌。此外,前辈们所宣讲和实行的数学表达方式,(至少可以说)给了工作的技术面以优先,而不大鼓励那些停留在“动机”上的“离题话”;甚至对那些仿佛要从雾里唤出某个也许富有启发性的形象或眼光的离题话也不鼓励——因为它们还没有化身为木头的、石头的或纯硬水泥的、可触摸的建筑,更像是梦的碎片,而不是工匠专注而尽责的活计。\origpage{P18}

从量上说,我在这些高强度多产的年份里的工作,具体表现为大约一万二千页的发表物——论文、专著或讨论班讲义\fnote{19}{}{从1960年代起,其中一部分发表物是与同行(尤其是 J. 迪厄多内)和学生合作写成的。}——以及数以百计(如果不是数以千计)的新概念,它们连同我提取它们时给的名字一起,进入了共同的遗产\fnote{20}{}{这些概念中最重要的那些,《主题概述》以及随附的《历史评注》都作了回顾,两者将收入《反思》第4卷。有些名字是朋友或学生向我建议的,比如“光滑态射”这个说法(J. 迪厄多内),或“site、champ、gerbe、lien”这一整套,那是让·吉罗(Jean Giraud)的博士论文里发展出来的。}。在数学史上,我想我大概是在我们这门科学里引入新概念最多的人,同时也正因如此,不得不发明最多新名字的人——为了把这些概念表达得细致,并尽我所能地富有提示性。\origpage{P19}

这些纯“量”的指标,当然只给出对我作品的一种极其粗略的把握,错过了真正构成其灵魂、生命与力量的东西。正如我刚才写的,我在数学里贡献的最好的东西,是我先得以窥见、随后耐心提取并或多或少加以发展的新“视角”。像我刚谈到的那些概念一样,这些新视角进入极其多样的情境,它们本身也几乎数不清。

然而有些视角比另一些更广阔,单凭它们就引发并包揽大量局部视角,涉及大量各不相同的具体情境。这样一个视角,也可以恰当地称为一个“大想法”。凭它自身的丰产性,这样一个想法生出成群的后代,这些想法都继承它的丰产性,但大多数(如果不是全部)比母想法范围更小。

至于表达一个大想法,也就是“说出”它,这往往是几乎和它被构想出来、和在构想者心中缓慢孕育一样精细的事——或者更确切地说,这费力的孕育与成形工作,恰恰就是“表达”该想法的那项工作:耐心地、日复一日地把它从诞生时环绕它的雾幕中提取出来,逐渐给它一个可触的形态,成为一幅在星期、月份和年份中不断丰富、稳固、精细起来的图画\origpage{P20}。仅仅给这想法命名,用某个醒目的说法,或用多少带点技术性的关键词,也许只要几行、甚至几页——但很少有人在尚未相当熟悉它的情况下,能听懂这个“名字”,并在其中认出一张面孔。而当这想法到了完全成熟的时候,也许一百页就足以表达它,让孕育它的那位工匠完全满意——也可能一万页经过长期打磨和掂量的文字仍不够\fnote{21}{}{1970年离开数学舞台时,我关于概形这一中心主题的全部发表物(其中不少是合作),加起来该有一万来页。然而这在我眼前看到的那个关于概形的宏大纲领里,只占不大的一部分。这纲领在我离开后便被无限期搁置了,尽管此前已经发展并发表、可供所有人使用的东西,除极少数例外,一开头就进入了共同的遗产,成为大家日常使用、视为“人所共知”的概念与结果。\par 我离开时已完成的那部分纲领,关于概形主题及其延伸与分枝,单是它本身就代表了数学史上所完成的最庞大的基础工作,而且肯定也是科学史上最庞大的之一。}。

而无论在哪种情形下,在那些为了把想法据为己有而了解了这项终于呈现出想法的蓬勃生长的工作的人当中——它像一片在林间荒地上长起来的宽阔乔木林——很有理由打赌:许多人会清楚看见所有这些挺拔有力的树,并且各有所用(有人为了爬上去,有人为了取梁木和木板,还有人为了在自己壁炉里烧旺火……),但很少有人懂得看见森林……

\section{(8) 眼光——或十二个主题的和谐}

也许可以说,“大想法”就是那样一种视角:它不仅显得新而丰产,而且把一个新的、广阔的主题引入科学,而这主题正是它的化身。而任何一门科学,当我们不把它当作权力与支配的工具,而当作我们这个物种跨越时代的认识冒险来听时,它不外就是这种和谐——在不同时代或广或狭、或贫或富——它在世代与世纪中展开,靠的是所有先后出现的主题之间那种精细的对位,那些主题像被从虚无中唤来,加入其中并在其中交织。\origpage{P21}

在我于数学中提取出的众多新视角里,回头看去,有十二个我会称之为“大想法”\fnote{22}{}{以下是这份十二个主导想法、即我作品的“主主题”的清单(按出现的时间顺序),供好奇的数学读者一阅。\par 1. 拓扑张量积与核空间。\par 2. “连续”与“离散”的对偶(导出范畴、“六种运算”)。\par 3. 黎曼–罗赫–格罗滕迪克瑜伽($K$-理论,与相交理论的关系)。\par 4. 概形。\par 5. Topos。\par 6. 平展上同调与 $\ell$-进上同调。\par 7. Motif 与 Motivic 伽罗瓦群(Grothendieck 的 $\otimes$-范畴)。\par 8. 晶体与晶体上同调,“德拉姆系数”“霍奇系数”的瑜伽……\par 9. “拓扑代数”:$\infty$-栈、derivator;topos 的上同调形式主义,作为新的同伦代数的灵感。\par 10. 温和拓扑学。\par 11. 远阿贝尔代数几何的瑜伽,伽罗瓦–泰希米勒理论。\par 12. 正多面体与各种正则构型的“概形”或“算术”视角。\par 除这些主题中的第一个(其中重要的一支属于我1953年的博士论文,并在1950至1955年间我的泛函分析时期得到发展)之外,其余十一个都是在我1955年以后的几何学家时期提取出来的。}。看我的数学家作品,“感受”它,至少是稍稍看见并“感受”其中某些想法,以及它们引入的那些构成作品的经纬与灵魂的大主题。

% ---------------- p.45 ----------------
事所必然,这些想法中有些比另一些“更大”(而那些也就因此“更小”!)。换句话说,在这些新主题里,有些比另一些更广,有些更深地潜入数学事物之谜的核心\fnote{23}{}{在这些主题中,就意义之广而言,我觉得最广的是 topos,它提供了把代数几何、拓扑学与算术综合起来的想法。就已产生的发展之规模而言,最广的是概形主题。(关于这一点,见第20页页底注(*)。)正是它为所论主题中的另外八个(即除第1、5、10之外的所有主题)提供了“典范的”框架,同时为代数几何以及代数—几何语言的自下而上的更新提供了中心概念。\par 在另一端,十二个主题中的第一个和最后一个,在我看来比其余的规模更有限。不过,就最后一个而言,它把一种新视野引入正多面体与正则构型这个极古老的主题,我怀疑一个全身心投入其中的数学家,一生也不足以把它穷尽。至于所有这些主题中的第一个,即拓扑张量积,它扮演的更多是一件现成新工具的角色,而不是后人发展的灵感来源。不过,直到最近这些年,我仍偶尔收到多少较近的工作的零星回响,那些工作(二三十年后)解决了我留下的某些问题。\par 在我看来,这十二个主题中最深的是 Motif,以及与之紧密相连的远阿贝尔代数几何和伽罗瓦–泰希米勒瑜伽。\par 从工具经过我亲手调校、臻于完善、并在近二十年研究的各个“前沿地带”日常使用的力量看,我觉得最值得注目的是“概形”与“平展上同调、$\ell$-进上同调”这两块。对一个消息灵通的数学家来说,我想从现在起几乎不可能怀疑:概形工具,以及由它产生的 $\ell$-进上同调工具,属于本世纪少数几项重大成就之列,它们在最近几代人之中滋养并更新了我们的科学。}。有三个(在我看来并非最不重要的)只是在我离开数学界之后才出现的,至今仍处胚胎状态;“正式地说”,它们甚至并不存在,因为没有任何正式发表物可以充当它们的出生证明\fnote{24}{}{这三者得到哪怕最简略勾勒的唯一“半正式”文本,是1984年1月为申请 CNRS 脱产而写的《纲领概要》。这一文本(导论第3节“罗盘与行囊”里也谈到它)原则上将收入《反思》第4卷。}。在我离开之前出现的九个主题里,最后三个我留在蓬勃生长之中,今天仍处在童年状态,因为(我走后)没有充满爱的手去照料这些“孤儿”的所需,它们在一个敌对的世界里被弃之不顾\fnote{25}{}{在我不久离开之后,这三个孤儿便被无声无息地埋葬了;其中两个后来又被大吹大擂地掘出,却不提工匠的名字——一个在1981年,另一个(鉴于此举毫无瑕疵的成功)就在次年。}。至于另外六个主题,在我离开前的二十年里已完全成熟,可以说(有一两点保留\fnote{26}{}{这“除了极少例外”,主要指格罗滕迪克式的对偶瑜伽(导出范畴与六种运算),以及 topos 的瑜伽。这一点将在《收获与播种》第二部和第四部(《葬礼(一)》与《葬礼(三)》)里详加讨论。})它们那时已经进入共同的遗产:尤其在几何人中间,今天“人人”都在整天、随时吟诵它们,甚至不再知道自己是在吟诵(就像茹尔丹先生说着散文)。它们是我们呼吸的空气的一部分,当我们“做几何”,或当我们做多少带点“几何性”的算术、代数或分析时。\origpage{P22}

我作品这十二个大主题绝不是彼此孤立的。在我眼里,它们属于一种精神与话题的统一,这统一像一个共同而持久的低音,贯穿我全部“写出的”和“未写出的”作品。写这些行的时候,我仿佛又听见了同一个低音——像一个召唤!——贯穿那三年“无偿的”、顽强的、孤独的工作,那时候我还不曾操心世上除我之外是否还有别的数学家,因为我当时完全被那召唤我之物的魅力攫住了……\origpage{P23}

这种统一,并不只是同一个工匠在他手下出来的作品上留下的印记。这些主题由无数既精细又显然的纽带彼此相连,就像那些各各清晰可辨的不同主题之间的相连:它们
% ---------------- p.46 ----------------
在同一片广阔的对位中展开、交缠——在一个把它们聚合、推向前去、并赋予每一个以意义的和谐之中,每一个的意义、运动与圆满都有所有其他者参与。每一个局部主题似乎都从这更广的和谐中生出,又随时从中重生,远不是这和谐显得像那些先于它存在的构成性主题的“总和”或“结果”。说实话,我无法摆脱这种感觉(大概很可笑……):在某种意义上,正是这尚未出现、却无疑早已确确实实“存在”于尚未诞生之物的幽暗怀抱中某处的和谐——正是它先后唤起了这些只有通过它才获得全部意义的主题,也正是它在我刚出少年、那些炽烈孤独的年月里,已经用低而急迫的声音呼唤我……

无论如何,我作品的这十二个主主题确实都像出于一种秘密的预定,共同汇成同一部交响曲——或者换一个比喻,它们体现为同样多的不同“视角”,全都汇入同一个广阔的眼光。

这眼光只是到1957、58年前后才开始从雾中浮现、显出可辨认的轮廓——那是紧张孕育的年份\fnote{27}{}{1957年正是我被引向提取“黎曼–罗赫”主题(格罗滕迪克版)的那一年——这个主题一夜之间把我捧成了“大明星”。这一年也是我母亲去世之年,因此也是我生命中一道重要分界线之年。那是我一生中创造力最旺盛的年份之一,而且不限于数学层面。那时,全部精力投入数学工作已有十二年。就在那一年,一种感觉浮现出来:关于什么是数学工作,我大致已经“转完了一圈”,也许现在是时候投入别的事情了。那显然是一种内在更新的需要,一生中第一次浮出水面。那时我想过要当作家,并且有好几个月停止了全部数学活动。最后我决定,至少还要把手上已有的数学工作写成白纸黑字,也就是几个月的事,顶多一年……\par 无疑,那一跃的时机还没有成熟。总之,一旦重新拾起数学工作,便是它重新攫住了我。此后又十二年,它再没有放开我!\par 这段间歇之后的一年(1958),也许是我数学生涯中最丰产的一年。新几何的两个中心主题正是在这一年绽放:概形理论的强力启动(这是我当年夏天在爱丁堡国际数学家大会上报告的主题),以及“site”概念的出现——它是关键的 topos 概念的技术性临时版本。隔了将近三十年回头看,我现在可以说:新几何的眼光真正诞生,就在这一年;它紧随这门几何的两件主要工具之后:概形(它代表旧“代数簇”概念的一次变形),以及 topos(它代表空间概念的一次更深层的变形)。}。也许有点奇怪:这眼光对我如此切近,如此“显然”,以至于直到一年前\fnote{28}{}{我头一次想到给这眼光一个名字,是在1984年12月4日的反思里,在《阴,仆役(二)——或慷慨》一注(ReS III,第637页)的第1361号子注中。},我都没想过给它一个名字。(而我的一大癖好,恰恰是不停地给向我显现的事物命名,作为把握它们的第一步……)确实,我说不出某一个特定时刻,可以被经历为这眼光出现的时刻,或事后可以认作那样的时刻。一种新眼光是那样广阔,它的出现无疑不能定在某一个特定时刻,而必须在漫长的年月里(如果不是在世代之中)逐步渗入并占有那个或那些细看与凝视它的人;仿佛在那些注定要被逐步取代的熟悉的眼睛后面,必须辛苦地长出新眼睛来。而这眼光也太大,谈不上像把握路上拐角处出现的第一个概念那样去“把握”它。所以,最终不必惊讶:给这样广阔、又如此切近而弥散的东西命名的念头,只是事后、只是当这东西达到完全成熟之后才出现。\origpage{P24}

说实话,直到两年前,我与数学的关系(除了教书的任务之外)仅限于做数学——跟从一股不断把我向前拉的冲动,进入一个不断吸引我的“未知”。我从未想过在这股冲劲中停下来,哪怕只停一瞬\origpage{P25}
% ---------------- p.47 ----------------
,回头看看也许已画出的走过之路,甚至为一部已过去的作品定位。(无论是把它定位在我的生命里,当作一件仍有深的、长期未被察觉的纽带把我与它相连的东西;还是把它定位在“数学”这场集体冒险里。)

还有一件怪事:要让我终于“放下”并重新认识这部半被遗忘的作品,或者仅仅想到给那曾是它灵魂的眼光一个名字,竟需要我突然面对一场规模巨大的“葬礼”的现实:那场以沉默和嘲弄进行的葬礼,埋葬的是那眼光,以及生出那眼光的那位工匠……

\section{(9) 形式与结构——或事物的道路}

不知不觉,这篇“前言”一针引一线,最后成了对我作品的一次正式介绍,主要是给非数学读者看的。已经陷得太深,退不回去了,我只剩下把“介绍”讲完!我想尽力至少说几句,谈谈那些我在前面篇幅里晃动过的奇妙“大想法”(或“主主题”)的实质,以及那道著名的“眼光”的本性——那些主导想法据说要汇入其中。由于无法求助于多少带点技术性的语言,我大概只能传递一个极其模糊的形象(如果真有什么东西能“传过去”的话……)\fnote{29}{}{这个形象不得不“模糊”,丝毫也不妨碍它是忠实的,也不妨碍它确实再现了所看之物的本质的某种东西(在此即我的作品)。反过来,一个形象无论多么清晰,也可能完全走样,而且只包含次要的东西,完全错过本质。所以,如果你“接住”了我关于自己作品所要说的东西(那肯定我心中形象的一部分会真的“传过去”),你就可以得意地说,你对我的作品中最本质的东西,把握得比任何一位博学的同行也许都更好!}。

传统上,人们区分宇宙之物的三类“性质”或“方面”作为数学思考的对象:数\fnote{30}{}{这里指的是所谓“自然数”0、1、2、3等等,或者(至多)指那些借它们、通过初等性质的操作表达出来的数(如分数)。这些数不像“实数”那样,可以用来度量一个可连续变化的量,比如直线上、平面上或空间中两个变动点之间的距离。}、量、形。也可以把它们叫做事物的“算术”方面、“度量”(或“分析”)方面、“几何”方面。在数学所研究的大多数情境里,这三个方面同时在场并紧密互动。然而,最常见的是其中一个明显占主导。在我看来,对大多数数学家来说(对认识他们或了解他们作品的人来说),他们的基本气质相当清楚:是“算术家”、“分析家”还是“几何家”——即便他们本领很多,在各种音区和调门里都工作过。\origpage{P26}

我最初那些孤独的思考,关于测度与积分理论,毫无歧义地属于“量”或“分析”一类。我引入数学的第一个新主题也是如此(它在我看来比另外十一个范围要小)。我之从“分析”这个“侧门”进入数学,在我看来并非由于我的特殊气质,而是由于可以叫做“偶然境遇”的东西:对于我这颗爱好普遍性与严格性的头脑来说,中学和大学向我提供的教学里最大的空缺,恰好在事物的“度量”或“分析”方面。

1955年是我数学工作的一个关键转折:从“分析”
% ---------------- p.48 ----------------
转向“几何”。我还记得那种震撼的印象(当然完全是主观的):仿佛我离开了干瘦而硌人的草原,忽然来到某种“应许之地”,那里丰饶的财富无穷无尽地增殖,凡是手愿意落下去采摘或挖掘的地方都是如此……而这种压倒性的丰饶印象,超出了任何尺度\fnote{31}{}{我用“压倒性的,超出任何尺度”这个词语组合,勉强译出德文的“überwältigend”及其英文对应词“overwhelming”。上一句里那个(并不恰当的)“震撼的印象”,也应当带着这层意味来理解:当面对一种非凡的壮丽、伟大或美时,我们心中激起的印象与情感突然把我们淹没,以致任何表达我们所感的意愿都仿佛预先被摧毁了。},在此后这些年里只得到确认和加深,直到今天。\origpage{P27}

这就是说,如果在数学里有一件东西(大概从来如此)比任何别的更令我着迷,那不是“数”,也不是“量”,而始终是形。而在形为了向我们显现而选取的千零一张面孔中,比任何别的更令我着迷、并继续令我着迷的,是数学事物中隐藏的结构。

一件事物的结构绝不是我们可以“发明”的东西。我们只能耐心地把它揭示出来,谦卑地去认识它,“发现”它。这项工作里若有创造性,我们若有时做着铁匠或不倦的建造者的活,那绝不是为了“塑造”或“建造”什么“结构”。那些结构根本不曾等我们才存在,才恰好是它们所是的样子!而是为了尽可能忠实地表达我们正在发现和探测的这些事物,以及这个不肯轻易交付的结构——我们摸索着,用也许还结结巴巴的语言,试着把它圈定。于是我们不断地“发明”适合的语言,以越来越精细地表达数学事物的内在结构,并借助这语言,“建造”那些理应说明已被把握和看见之物的“理论”,一点一点、从头到尾。这里有一个持续不断的往返运动:在把握事物与表达被把握者之间,那语言在工作中、在当下需要的不断压力下变得精细并重新被创造。

读者大概已经猜到,这些“从头建造”的“理论”,也正是前面谈到的那些“漂亮的房子”:我们从先人那里继承的房子,以及我们在事物的召唤与倾听中亲手建造的房子。而我刚才说到建造者或铁匠的“创造性”(或想象力),还得补充一句:构成其灵魂与秘密神经的,绝不是那个说“我要这个,不要那个!”、并乐于凭自己心意决断的人的傲气;那就像一个糟糕的建筑师,还没看见、感受过一块地,还没探明它的可能与要求,脑子里就已经准备好了图纸。研究者的创造性与想象力的品质,在于他注意力的品质,在于他倾听事物之声。因为宇宙之物从不厌倦谈论自己、向那个有心去听的人显现。而最美的房子,显出工匠之爱的那一座,并不是比别人更大或更高的那座。美的房子,是忠实映照事物隐藏的结构与美的那座。\origpage{P28}

\section{(10) 新几何——或数与量的联姻}

可我又跑题了——我本打算谈主主题,它们在同一位母眼光中结合起来,像许多条河流回归它们为其子嗣的大海……

% ---------------- p.49 ----------------
这广阔的统一性眼光可以描述为一种新几何。据说,这就是上个世纪克罗内克所梦想的那一种\fnote{32}{}{关于这个“克罗内克的梦”,我只是听人说起:有人(也许正是约翰·泰特)对我说,我正在实现那个梦。在我从前辈那里接受的教育里,历史性的参考极少见;哺育我的并不是阅读多少古代的、甚至同时代的作者,而主要是与其他数学家——首先是我的前辈们——的口头或通信交流。1958年概形理论突然而强劲的启动,其主要的、也许是唯一的外来灵感,是塞尔那篇以 FAC(《代数凝聚层》)这个缩写广为人知的文章,它发表于此前几年。除此之外,理论后来发展中的主要灵感都是从自身产生的,并在年复一年中自我更新,靠的只是内部简洁性与一致性的要求,以及一种努力:在新语境中说明代数几何里“人所共知”的东西(那些东西在我手中变形的同时被我逐步消化),以及这个“已知”让我预感到的东西。}。但现实(一个大胆的梦有时让人预感到或窥见它,并鼓励我们去发现它……)每次都在丰富与共鸣上超过最大胆、最深的梦。可以肯定,对这新几何的不止一个部分(如果不是全部),在它出现的前一天,没有人会想到——工匠本人也不比别人更早想到。

可以说,“数”适合把握“不连续”或“离散”的聚集体的结构:那些往往有限的系统,由可以说是彼此孤立的“元素”或“对象”组成,没有某种从一个到另一个的“连续过渡”原则。“量”则相反,是典型地能够“连续变化”的性质;因此它适合把握连续的结构与现象:运动、空间、各种“簇”、力场等等。于是,算术显得(大致上)是离散结构的科学,分析则是连续结构的科学。

至于几何,可以说,它以现代意义上的科学的形式存在了两千多年,一直“骑”在这两类结构上,即“离散的”与“连续的”\fnote{33}{}{说实话,传统上处在几何家注意力中心的是“连续”方面,而“离散”性质,尤其是数值的和组合的性质,则被略而不谈,或草草打发。十年前我惊叹地发现二十面体组合理论的丰富,而这个主题在克莱因论二十面体那本经典著作里连碰都没有碰到(大概根本没看到)。我看到几何家对自发进入几何的离散结构这种(长达两千年的)忽视的另一个显著迹象:群(尤其对称群)的概念直到上世纪才出现,而且它最初还是(由埃瓦里斯特·伽罗瓦)在一个当时并不被视为属于“几何”的语境中引入的。确实,直到今天,仍有许多代数学家始终不明白:伽罗瓦理论本质上正是一种“几何”眼光,它更新了我们对所谓“算术”现象的理解……}。而且长期以来,两种本属不同种类的几何——一种离散、一种连续——之间并没有真正“离婚”。不如说,那是对同样的几何图形作研究时的两种不同视角:一个强调“离散”性质(特别是数值的和组合的性质),另一个强调“连续”性质(比如在周围空间中的位置,或以点的相互距离来度量的“量”等等)。\origpage{P29}

直到上世纪末,才出现了一场离婚,随着人们有时称为“抽象(代数)几何”的东西的出现与发展。大致说来,它在于为每个素数 $p$ 引入一种“特征 $p$ 的(代数)几何”,照着前几个世纪传下来的(代数)几何的(连续)模型来仿制,然而所处的语境却显得不可还原地“不连续”“离散”。这些新的几何对象自本世纪初以来地位日益重要,尤其是由于它们与算术——离散结构的典型科学——的紧密关系。这似乎是安德烈·韦伊(André Weil)\fnote{34}{}{安德烈·韦伊,旅居美国的法国数学家,是“布尔巴基小组”的“创始成员”之一;《收获与播种》第一部里会不少谈到这个小组(偶尔也会谈到韦伊本人)。}作品中的一条指导性想法,也许还是主要的力量线(在他的书面作品中,它或多或少停留在默认状态,这也是应该的):即“代数”几何,尤其是与不同素数相联系的“离散”几何,应当提供钥匙,让算术得到大范围更新。他就是本着这种精神,在1949年提取出著名的“韦伊猜想”的。

% ---------------- p.50 ----------------
说实话,这些猜想令人目瞪口呆:它们让人们对这些离散本性的新“簇”(或“空间”)隐约看见某些类型的构造与论证\fnote{35}{}{(供数学读者参考。)这里指的是与微分簇或复簇的上同调理论相联系的“构造与论证”,尤其是涉及莱夫谢茨不动点公式与霍奇理论的那些。}的可能性——而在此前,这些似乎只在分析家认为配得上这个名号的“空间”,即所谓“拓扑”空间(连续变化的概念在其中通行)的框架里才可设想。\origpage{P30}

可以认为,新几何首先是这两个世界之间的综合——它们此前相邻、紧密相依,却又彼此分隔:“算术”世界,其中住着那些没有连续性原则的(所谓)“空间”;以及连续量的世界,其中住着本来意义上的“空间”,它们对分析家的手段敞开,(正因如此)也被他接受为配得上在数学之城里有住所。在新眼光里,这两个昔日分隔的世界只成了一个。

这种“算术几何”(我提议这样称呼这种新几何)眼光的第一个胚胎,就在韦伊猜想里。在我某些主要主题\fnote{36}{}{指四个“中间”主题(第5至8号),即 topos、平展上同调与 $\ell$-进上同调、Motif,以及(较小程度上)晶体。这些主题是我在1958到1966年间先后提取出来的。}的发展中,从1958到1969这些年,这些猜想始终是我主要的灵感来源。而且,在我之前,一边是奥斯卡·扎里斯基(Oscar Zariski),另一边是让–皮埃尔·塞尔,已经为“抽象”代数几何那些无法无天的空间发展出某些“拓扑”方法,其灵感来自此前用于大家那些“正经”空间的方法\fnote{37}{}{(供数学读者参考。)在我看来,扎里斯基在这方面的主要贡献是引入“扎里斯基拓扑”(它后来成了塞尔在 FAC 中的一件本质工具),以及他的“连通性原则”和他所谓“全纯函数论”——后者在他手中变成了形式概形理论,以及形式与代数之间的“比较定理”(第二个灵感来源是塞尔那篇奠基性的 GAGA)。至于我在正文中提到的塞尔的贡献,当然首先是他把层(由让·勒雷在十几年前、在一个完全不同的语境中引入)的视角引入抽象代数几何,即在那篇已经提到的另一篇奠基性文章 FAC(《代数凝聚层》)中。\par 借这些“回顾”之光,如果我要举出新几何眼光最近的“祖先”,立刻想到的就是奥斯卡·扎里斯基、安德烈·韦伊、让·勒雷和让–皮埃尔·塞尔这几个名字。其中塞尔起了独特的作用,因为主要正是通过他,我才不仅知道了他自己的想法,也知道了扎里斯基、韦伊和勒雷那些在新几何的绽放与发展中起过作用的想法的。}。

当然,他们的想法在我构筑算术几何的最初几步中起了重要作用;不过说实话,更多是作为出发点和工具(为了一个广阔得多的语境的需要,我不得不或多或少把它们从头重做),而不是作为一路继续滋养我的梦与计划的灵感来源。无论如何,一开始就很清楚:即便重做过,这些工具也远远达不到所需——连朝着那些奇妙猜想走出最初几步都不够。\origpage{P31}

\section{(11) 魔法扇——或天真}

在新几何的启动与发展中,两个关键的力量线是概形与 topos。它们几乎同时出现,紧密共生
% ---------------- p.51 ----------------
\fnote{38}{}{这次启动发生在1958年,第23页的页底注里谈到它。“site”或“格罗滕迪克拓扑”的概念(topos 概念的临时版本)紧随着概形概念出现。反过来,正是它提供了“局部化”或“下降”的新语言,在概形主题与工具的发展中每一步都要用到。更为内蕴、更具几何性的 topos 概念,起初在随后几年里仍是隐性的,主要从1963年起随着平展上同调的发展而浮现,并逐步在我这里确立为最根本的概念。},在新几何引人注目的兴起中像是同一根驱动神经,而且从它们出现的那一年起就如此。为了结束这趟对我作品的巡览,我至少还得就这两个想法说几句。

概形这个概念,是把此前人们摆弄的那无穷系列的(代数)“簇”概念(每个素数一个这样的概念\fnote{39}{}{这一系列里还应当包括 $p=\infty$ 的情形,它对应“特征零”的代数簇。}……)统摄进唯一概念的想得到的最自然、最“显然”的办法。此外,同一个“概形”(或新版“簇”)对每个素数 $p$ 都产生一个完全确定的“特征 $p$ 的(代数)簇”。这些不同特征的簇的集合,于是可以被看成某种“簇的(无穷)扇”(每个特征一把)。而“概形”就是这把魔法扇,它把自己的各种可能特征下的“化身”或“转世”彼此连起来,像许多不同的“扇骨”。由此,它提供了一个有效的“过渡原则”,把那些此前显得多少彼此孤立、互不相通的几何所属的“簇”彼此连起来。如今它们被统摄在一门共同的“几何”里,并由它相连。可以把它叫做概形几何,它是那种“算术几何”的最初草图,而后者将在随后几年里在其中充分展开。\origpage{P32}

概形这个想法本身天真地简单——如此简单、如此卑微,以致在我之前没人想到要弯得这么低。说实话,甚至“蠢”得可以:以致好些年里,尽管事情明摆着,对我许多博学的同行来说,这实在“不严肃”!而我也确实需要好几个月紧凑而孤独的工作,才能在自己的角落里说服自己:“它行得通”——我那股不可救药的幼稚使我执意要试验的新语言,那个蠢得可以的语言,确实适合以新的光亮和新的精细,并且从此在一个共同的框架里,把握那些与前此“特征 $p$ 的几何”相连的最早的几何直觉中的某些。这类练习,任何“消息灵通”的人都会预先判定为愚蠢而无望;而在我的所有同行和朋友中,大概只有我一个人会想到把它放进脑子,甚至(被一个秘密的魔鬼驱使……)不顾所有人反对把它做到底!

我并没有让自己被周围那些关于什么“严肃”、什么不严肃的立法共识分散注意力,而是像过去一样,单纯地信任事物卑微的声音,信任我身上那个懂得倾听的东西。报偿是即时的,而且超出一切期待。在这几个月里,我甚至没有“特意”,就摸到了强大而出人意料的工具。它们让我不仅能以更透彻的光重新找到(像玩一样)那些公认难懂的旧结果并超越它们,而且终于能够着手并解决一些“特征 $p$ 的几何”问题——此前用当时一切已知手段都显得无法企及\fnote{40}{}{概形理论这次“强力启动”的报道,就是我1958年在爱丁堡国际数学家大会上的报告。在我看来,这篇报告是概形视角最好的导论之一,(也许)能让一位几何读者有动力去勉强啃一啃那部(后来的)厚重论著《代数几何原理》——它详详细细地(一个技术细节都不放过)阐述了代数几何的新基础与新技巧。}。\origpage{P33}

% ---------------- p.52 ----------------
在我们对宇宙之物的认识里(无论数学的或别的),我们身上那革新的力量不外是天真。那是我们出生时人人分得的原初天真,它安息在每个人心里,常常成为我们蔑视的对象,也成为我们最隐秘的恐惧的对象。唯有它把谦卑与大胆结合起来,让我们深入事物之心,也让我们允许事物深入我们、浸透我们。

这种力量绝不是“天赋”异禀的特权——比如说有一种超常的脑力,能灵巧轻便地消化和摆弄大量已知的事实、想法和技术。这样的天赋当然可贵,对于像我这样出生时没有被“超乎一切尺度地”赋予它的人,当然值得羡慕。

然而,不是这些天赋,也不是最炽烈的野心配上无懈可击的意志,能让人越过那些把我们的宇宙关起来的“看不见而专横的圈子”。唯有天真能越过它们,不知不觉也不在意,在那些我们独自倾听事物、全神贯注于一场儿童游戏的时刻……

\section{(12) 拓扑学——或雾中的丈量}

我们刚看到,“概形”这个革新性的想法,让人能把与不同素数(或不同“特征”)相联系的不同“几何”彼此连起来。然而这些几何各自本质上仍是“离散”或“不连续”的,与过去几个世纪传下来的(可上溯到欧几里得的)传统几何相对照。扎里斯基和塞尔引入的新想法,在一定程度上为这些几何恢复了某种连续性的“维度”,它随即被刚出现、旨在把它们统一起来的“概形几何”继承。但就那些(韦伊的)“奇妙猜想”而言,还差得远。从这个角度看,这些“扎里斯基拓扑”粗疏到这种地步,几乎就像还停留在“离散聚集体”的阶段。显然缺少的是某种新原则,能把那些几何对象(或“簇”,或“概形”)与通常的、或“正经”的(拓扑)“空间”连起来;所谓正经的空间,姑且说,是那些“点”显得彼此清楚分开的空间,而在扎里斯基引入的那些无法无天的空间里,点却有一种讨厌的倾向,彼此黏成一团……\origpage{P34}

确实,只有这样一个“新原则”的出现,而且非它不可,才能让那场“数与量的联姻”,或者说“不连续几何”与“连续几何”的联姻得以完成——韦伊猜想里已经有了对它的第一次预感。

“空间”概念大概是数学里最古老的概念之一。它在我们对世界的“几何”把握中如此基本,以致两千多年来它或多或少一直是默认的。只是在上一个世纪里,这个概念才逐步摆脱直接知觉(对我们周围那唯一同一个“空间”的知觉)及其传统理论化(“欧几里得式”)的专制控制,获得自身的自主与动力。今天,它属于数学中最普遍、最常用的少数概念之列,无疑为无一例外的所有数学家所熟悉。它也是再千变万化不过的概念,有千百张面孔,取决于人们给这些空间加上什么类型的结构:从最丰富的(如那些可敬的“欧几里得”结构,或“仿射”与“射影”结构,或同名“簇”的“代数”结构,后者推广并柔化了前者),直到最贫乏的:那些任何“定量”信息要素似乎都一去不返地消失了
% ---------------- p.53 ----------------
,只剩“邻近”或“极限”概念之定性精髓\fnote{41}{}{说到“极限”概念,我在这里想的尤其是“取极限”那个概念,而不是(对非数学人士更熟悉的)“边界”那个概念。},以及形式直觉最难以捉摸的版本(所谓“拓扑的”)。在这些概念中,最贫乏的一个,也是过去半个世纪里一直充当某种广阔的公共概念怀抱、把其余一切统摄于其中的那个,就是拓扑空间。对这些空间的研究构成几何中最迷人、最有生命力的分支之一:拓扑学。\origpage{P35}

由一个(所谓“拓扑”)“空间”所体现的这种“纯品质”结构,乍看之下无论多么难以捉摸——因为一切定量性质的数据(特别是两点之间的距离)都缺席,使我们无法挂靠任何关于“大”“小”的熟悉直觉——上个世纪里人们毕竟还是把这些空间精细地圈在了一套仔细“量身剪裁”的语言的紧密而柔韧的网眼里。更妙的是,人们还从头发明并制造出各种“米尺”或“量绳”,无论如何也要给那些触手般伸展的“空间”挂上各种“度量”(叫做“拓扑不变量”)——那些空间像抓不住的雾,似乎躲开一切测量的尝试。确实,这些不变量中大多数、而且是最本质的那些,其本性比单纯的“数”或“量”更精细——它们本身就是多少有些精巧的数学结构,以多少复杂些的构造挂到所考虑的空间上。这些不变量中最古老、最关键之一,上个世纪已经(由意大利数学家贝蒂)引入,是由与该空间相联系的各个所谓“上同调”“群”(或“空间”)构成的\fnote{42}{}{说实话,贝蒂引入的不变量是同调不变量。上同调是它的一种大致等价、“对偶”的版本,出现得晚得多。这个方面压过了最初的“同调”方面,主要(无疑)是由于让·勒雷引入了层的视角,下面会谈到。从技术角度看,可以说我大部分几何学家的工作,都在于提取并或多或少地发展那些缺失的上同调理论,为各种空间与簇、尤其是为“代数簇”和概形而设。一路之上,我还被引向用上同调的语言重新解释传统的同调不变量,并由此让它们在一种全新的光下显现。\par 拓扑学家还引入了许多别的“拓扑不变量”,用来把握拓扑空间的这类或那类性质。除了空间的“维数”和(上)同调不变量,最早的其他不变量是“同伦群”。1957年我又引入了另一个,即所谓“格罗滕迪克”群 $K(X)$,它随即大行其道,其重要性(无论在拓扑学还是在算术中)至今不断得到确认。\par 一大批新不变量,本性比目前已知和使用的更精细,但我觉得是根本性的,都写进了我的“温和拓扑学”纲领(其极简的概述见《纲领概要》,将刊于《反思》第4卷)。这一纲领基于“温和理论”或“温和空间”的概念,它有点像 topos 那样,构成空间概念的(第二次)“变形”。它比后者显然得多(在我看来),也没有那么深。我预计它对“严格意义上的”拓扑学的直接影响会明显更有冲击力,并且会彻底改变几何拓扑学家的“手艺”,办法是深刻改造他所工作的概念语境。(正如代数几何里引入概形视角时所发生的那样。)此外,我把自己的《纲领概要》寄给了好几位旧友和著名的拓扑学家,但看来它并没有引起他们中任何一位的兴趣……}。\origpage{P36}

正是它们(确实,主要是“在字里行间”)出现在韦伊猜想里,构成其深刻的“存在理由”,并使它们(至少对我来说,被塞尔的解释“带进了门”)具有全部意义。但要把这样的不变量挂到这些猜想里出现的“抽象”代数簇上,以回应为那一事业的需要而提出的非常精确的要求——这还只是一个希望。我怀疑,除塞尔和我自己之外,没有别人(甚至,尤其是,安德烈·韦伊本人!\fnote{43}{}{说来矛盾,韦伊对上同调形式主义有一种顽固的、看来是发自内脏的“阻塞”——而这恰恰主要是他那些著名的猜想,从1955年起启发了代数几何中那些宏大上同调理论的发展(塞尔以其已在前面的页底注中提到过的奠基性文章 FAC 开了第一炮)。\par 在我看来,这种“阻塞”在韦伊那里属于一种普遍的厌恶:厌恶一切“大堆头”,厌恶一切近似于形式主义(当它无法用几页纸概括时)或多少有点复杂交错的“构造”的东西。他确实一点也不像“建造者”;显然,他在三十年代是被迫地动手发展代数几何最初的“抽象”基础,而那些基础(鉴于他这种性情)对使用者来说实在成了一张普罗克鲁斯忒斯之床。\par 我不知道他是否怪我得太过分,怪我投入去建造那些宽敞的住宅,使克罗内克的梦和他的梦得以化身为一套精细而有效的语言与工具。总之,他从未就他看见我从事的工作、或已经完成的工作,对我说过一句话的评论。我把《收获与播种》寄给他三个多月,还附上我亲手写的热情题献,也没有得到任何回音。})真的相信……\origpage{P37}
